% =====================================================================
\section{Discussion}
\label{sec:discussion}
% =====================================================================

\subsection{What the feedback buys, and what it does not}
\label{sec:discussion-content}

The confirmatory result is a null, and a ceiling bounds what it can mean.
Measured-fact feedback and a contentless rejection repaired injected defects
equally well, to within about 4 percentage points. Both sat near the capability
ceiling, which \arm{A0} reached in every trial. For defects of this kind, on
this model, it does not matter what the supervisor says, provided the model is
told that something is wrong and can see what it produced. Restating the
stage's obligations at retry time did not matter either. That excludes the
explanation \citet{huang2024cannot} and \citet{kamoi2024} warn of, in which
measured feedback wins by teaching the task. The C1 gap and the C0 gap were
identical.

The mechanism suggests why. We offer it as an interpretation, not a tested
claim. Derivability works in both directions. The supervisor can detect an
injected defect because the defect is inconsistent with upstream state that the
pipeline holds. That same upstream state is in the retried stage's prompt. An
evidence identifier that matches none of the reviews the stage was given, or a
recommendation unrelated to any pattern it was given, is as visible to a capable
model holding its own output and its inputs as it is to the checker. The
measured fact repeats what the model can already see. On this reading the
copy-and-correct pattern is to be expected: the model keeps what is consistent
with its inputs and reverts what is not, and the result is often its original
output byte for byte.

The result therefore sides with the accounts that locate the active ingredient
in the signal and in what the model can see
\citep{chen2024selfdebug,chen2026rolerelabel,tsui2026scbench}. It does not side
with those that locate it in feedback content~\citep{olausson2024selfrepair},
but only for defects this easy to undo. The core says nothing about defects the
model cannot see or cannot fix. The one cell with headroom, the
production-like \op{O1} at 84.7\%, gives the same answer: \arm{A1} and \arm{A2}
were identical. Its interval of $\pm12.5$\pp{} is too wide to exclude a
moderate effect.

The second model supplies the headroom that the core lacks, and there a gap
opens: measured facts led by +6.8\pp{} \ci{+0.5}{+13.3}, mostly at the
reasoning stage (Section~\ref{sec:results-rq6}). The gap is exploratory and
fragile ($p=0.052$, and the first repetition alone includes zero), but its form
fits the mechanism above. Without the measured fact, the second model left the
defect in place after a quarter of its retries; with it, after fewer than one
in ten. The fact supplied the location of the inconsistency, which the hosted
model found unaided. Models that cannot find their reasoning errors can
correct many of them once told where they are~\citep{tyen2024}, and the second
model behaved that way. Location is not always enough: on \op{O1}, where the
fix is to recover an exact identifier, naming it helped neither model. On this
evidence the null is a property of a capable model at its ceiling, not of the
method. We read the gap as a hypothesis: feedback content matters to the
extent that the model cannot find the defect by itself.

\subsection{Placement is the design lever the production retry lacks}
\label{sec:discussion-placement}

The production retry does not show the model its rejected output; it runs at
\arm{P0}. Showing the output (\arm{P1}) left repair statistically unchanged
($+3.7$\pp{} \ci{-1.5}{+9.6}). It cut collateral damage from 38.4\% of retries
to 1.1\%, and it avoided a schema-correction call that every report-stage retry
at \arm{P0} needed. A model that cannot see its output regenerates all of it,
and regeneration churns content that was correct. The same comparison favours
showing the artifact at equal budget. One retry at \arm{P1} repaired 97.0\%
of defects, at about a quarter of the token cost of three fresh draws, which
repaired 93.3\%. These results are exploratory and come from one model. But they
point to a change that costs nothing to make and is invisible to every
repair-rate metric: return the rejected output with the rejection. The second
model did not test placement, and below the ceiling wording mattered too
(Section~\ref{sec:discussion-content}). But the production feedback already
carries the measured facts; showing the output is the change that remains.

\subsection{A supervisor's precision bounds the value of its retries}
\label{sec:discussion-precision}

On real, production-caught defects the production retry was accepted 28.2\% of
the time. The obvious reading is that real defects are harder to fix. The data
support a different one. Most real defects were checker false positives: the
pre-study audit found none of them to be an error to a human rater, and the
calibration measured the trace check firing on 5.0\% of grounded
recommendations. The retry cannot repair content that is already correct. It
regenerates the content in other words, and the same lexical check rejects it
again. The rate also shows in clean runs: 16\% of clean pipelines in the
downstream pass were rejected, every one by the trace check. A false positive
costs a model call, and at the retry cap it becomes a halt or a warning on a
run that had nothing wrong with it. For the engineer of such a supervisor, the
precision of each check matters more than the wording of the feedback. On our
data, a content-word matcher fired on no grounded issue and on every foreign
one. We kept the production supervisor unchanged, because it is the system
under study. Replacing its matcher is the obvious follow-up, and we have not
evaluated it.

The same result is a caution for fault-injection studies of LLM systems in
general. The injected defects coupled to the real ones at the level of check
class: every real defect fires a class that some operator targets. Their repair
rates were nonetheless 98.5\% and 28.2\%. For this system, the question
\citet{just2014mutants} raised for program repair has a negative answer.
A repair rate measured on injected defects is not an estimate of the rate on
the defects the system produces, and coupling by class does not make it one. A
study that reports only injected-defect repair should not describe it as a
production repair rate. Ours does not, and the replay is what entitles us to
say so.

\subsection{Local checks verify fidelity, not correctness}
\label{sec:discussion-local}

An unrepaired defect passed the downstream checks about as often as a clean
run, and the report-containment check \emph{rewarded} carrying it forward. The
reason is structural. A check that verifies each stage against its upstream
certifies consistency, and consistency with a defective upstream is still
consistency. Once a defect has been accepted where it arose, every later check
treats it as ground truth. This is a property of derivability-based
verification in general, not of our thresholds. The conference
version~\citep{hoang2026eas} noted that ``an accepted claim may inherit an
upstream error''; the downstream pass measures it. Two consequences follow.
First, such a supervisor needs its recall at the point of origin, because
downstream checks cannot make up for a miss
upstream. Second, repair is not restoration. Repaired runs removed the planted
content but still diverged from the clean report beyond the noise floor,
because a repaired stage is a regenerated stage. Together with the loose
coupling between defect-level and pipeline-level outcomes, this is the pattern
\citet{cemri2025mast} report for multi-agent failures. A supervisor that
improves the defect-level rate need not move the run-level outcome, and
evaluations should report both.

\subsection{What the overwrite bought}
\label{sec:discussion-overwrite}

The retry repaired every divergent frequency it was given, without collateral
damage. The recompute-and-overwrite design was therefore not needed for
correctness. It is still the better design, for reasons the counterfactual makes
precise: it is exact by construction, costs no call, cannot introduce churn, and
does not depend on the model's cooperation. The principle that a quantity the
checker can compute should not be asked of the model stands. It stands on cost
and determinism, not on any failure of the alternative.

\subsection{Revisions to the conference version}
\label{sec:discussion-reconcile}

The new results qualify four claims of~\citet{hoang2026eas}.
\begin{enumerate}
\item \emph{``At temperature zero an unchanged prompt returns an unchanged
  answer, so the content of that feedback, not the act of retrying, is what can
  change the outcome.''}
  Both halves need qualifying. On a hosted API, temperature~0 reproduced a
  trial's outcome class across sessions in 96\% to over 99\% of pairs, but its
  output byte for byte in only 76--83\%. Within a single session, 27\% of \arm{A1}
  cells still produced two distinct outputs across three repetitions. A local
  runtime serving one request at a time did better, but two of ten prompts sent
  twice still returned different text. And with the artifact in view, the
  content of the feedback did not change the outcome on the hosted model. On the
  second model it changed it modestly (Section~\ref{sec:results-rq6}).
\item \emph{Defect incidence (35\% of runs) and instability (20\%).} These
  rates count the supervisor's verdicts. Some of those verdicts are false
  positives of the trace check, which fires on 5.0\% of grounded
  recommendations. Some identifier flags are single-character corruptions where
  the real evidence is present. The rates should be read as rates of supervisor
  flags. The rate of genuinely ungrounded content is lower by an amount this
  study bounds but does not measure.
\item \emph{Recompute-and-overwrite} was asserted as a principle with no
  counterfactual. It is now measured: it was a cost decision, not a correctness
  decision (Section~\ref{sec:results-rq5}).
\item \emph{Retry repair} was unmeasured, because the enforcing arm produced one
  retry. The hosted model repairs injected defects at the ceiling under any
  feedback; the second model repairs 66--73\% of them. The
  real defects the supervisor caught are accepted after the production retry
  28.2\% of the time, for the reasons in
  Section~\ref{sec:discussion-precision}.
\end{enumerate}
% Author decision 3 (JOURNAL_PAPER_PLAN.md section 2), closed 2026-10-23:
% items 2 and 4 are journal findings, following the precedent set for the
% matcher defect on 2026-09-14, not a note to the conference chairs. If the
% FISAT camera-ready changes any of the four claims quoted above, revise this
% list to match the published version.

\subsection{Recommendations for building runtime supervisors}
\label{sec:discussion-recommendations}

Five practical lessons follow for engineers who put a deterministic checker in
the control path of an LLM pipeline. All are scoped to one pipeline. Lessons
1, 3 and 4 rest on the hosted model alone. The second model qualifies lesson 2
and extends lesson 5 to a local runtime.
\begin{enumerate}
\item \textbf{Return the rejected output with the rejection.} Repair stayed the
  same, collateral damage fell from 38\% to 1\%, and one retry repaired at
  least as often as three fresh draws, at under a quarter of their token cost.
\item \textbf{Invest in check precision before feedback wording.} Wording made
  no measurable difference on the hosted model. On the second model, naming
  the defect helped by an exploratory 6.8 points, apparently by showing the
  model where the defect was. Keep the measured facts, which cost nothing, but do not
  expect them to make up for an imprecise check. False positives made up most
  of what the production retry was asked to fix.
\item \textbf{Verify where defects arise.} Stage-local consistency checks
  certify fidelity to upstream, so they cannot catch a defect after a later
  stage has built on it.
\item \textbf{Compute what can be computed.} Overwriting a derivable statistic
  is exact, free and churn-free. A retry achieves the same correctness at a
  cost.
\item \textbf{Treat temperature~0 as non-deterministic, on a hosted API and on
  a local runtime alike, and check injected defects against real ones.} Match
  every comparison within a serving session. Replay real defects before reading
  an injected-defect repair rate as a production one.
\end{enumerate}
